\chapter{$E_1$-Chiral Algebras}
\label{ch:e1-chiral}

Volume~II (Part~III) developed the $E_1$-chiral theory as the ordered sector of the Swiss-cheese algebra.  This chapter reformulates it as the foundation for the $E_2$ extension.

% AP104: The E_1/ordered story is the PRIMITIVE layer.  The E_2 braided
% structure is obtained from it by the Drinfeld center construction.
% Quantum groups, Yangians, and braided tensor categories are natively E_1.
% The passage E_1 -> E_2 is the higher-categorical analogue of the
% averaging map av: g^{E_1} -> g^mod from Vol I.

\begin{remark}[$E_1$ primacy for CY quantum groups]
\label{princ:e1-primacy-vol3}
The $E_1$-chiral algebra (boundary) is the primitive object in this volume.
The $E_2$-chiral algebra (bulk) is obtained from it by the Drinfeld center
construction $\cZ(\Rep^{E_1}(A)) \simeq \Rep^{E_2}(\mathrm{Drin}(A))$.
Quantum groups, Yangians, and braided tensor categories are natively $E_1$
objects: the CoHA multiplication is ordered (short exact sequences have a
preferred direction), and the $R$-matrix arises only in the Drinfeld double.
The passage $E_1 \to E_2$ is the higher-categorical analogue of the
averaging map $\mathrm{av} \colon \fg^{E_1} \to \fg^{\mathrm{mod}}$
from Vol~I.
This averaging is \emph{lossy}: $\mathrm{av}(r(z)) = \kappa_{\mathrm{ch}}$ forgets the
$R$-matrix data.  The $E_1$-bar $B^{\mathrm{ord}}(A)$ retains
strictly more information than the $E_\infty$-bar $B^{\Sigma}(A)$.
\end{remark}

\section{$E_1$-algebras and associative factorization}
\label{sec:e1-algebras}

An $E_1$-algebra is an algebra over the little intervals operad.  In the chiral setting, an $E_1$-chiral algebra factorizes associatively along a real line $\mathbb{R}$ while maintaining holomorphic structure along a transverse curve.

\begin{definition}[$E_1$-chiral algebra]
\label{def:e1-chiral-algebra}
An \emph{$E_1$-chiral algebra} on $X \times \mathbb{R}$ is a factorization algebra $\cA$ on $\Ran(X) \times \Ran(\mathbb{R})$ such that:
\begin{enumerate}[label=(\roman*)]
  \item The $X$-direction factorization is holomorphic (chiral);
  \item The $\mathbb{R}$-direction factorization is $E_1$ (associative).
\end{enumerate}
The representation category $\Rep^{E_1}(\cA)$ is a monoidal category.
\end{definition}

\section{The bar complex as $E_1$-chiral algebra}
\label{sec:bar-e1}

\begin{proposition}
\label{prop:bar-e1-structure}
\ClaimStatusProvedElsewhere
For a chiral algebra $A$ on $X$, the bar complex $B(A)$ carries a natural $E_1$-chiral structure: the bar differential encodes the $X$-direction (holomorphic OPE residues) and the bar coproduct encodes the $\mathbb{R}$-direction (ordered deconcatenation).
\end{proposition}
% AP85: The factorization coproduct (Vol I, coshuffle/Sym^c) is DIFFERENT
% from the deconcatenation coproduct (Vol II, T^c).  The E_1-bar uses
% deconcatenation; the E_infty-bar uses coshuffle.  See AP82, AP85.

This is the Swiss-cheese algebra structure of Volume~II.  The $E_1$-bar
$B^{\mathrm{ord}}(A)$ is the ordered tensor coalgebra $(T^c(A[1]), d, \Delta_{\mathrm{decat}})$
with deconcatenation coproduct (AP82, AP85).  The symmetric bar
$B^{\Sigma}(A) = B^{\mathrm{ord}}(A) / S_n$ is the $S_n$-coinvariant quotient.
The ordered bar retains the full quantum group data (R-matrix,
Yangian structure) that the symmetric bar quotients out (AP65, AP97).

\section{Monoidal representation categories}
\label{sec:e1-monoidal-rep}

The representation category $\Rep^{E_1}(\cA)$ inherits a monoidal structure
from the ordered factorization along $\mathbb{R}$.  The tensor product
$V \otimes W$ of two $\cA$-modules is defined by the fusion of intervals:
$V$ is placed on $(-\infty, 0)$ and $W$ on $(0, \infty)$, and the
factorization structure combines them.  The ordering of the intervals makes
the tensor product associative but \emph{not} symmetric: the natural
transformation $V \otimes W \to W \otimes V$ does not exist in the
$E_1$-setting.  It exists only after passing to the Drinfeld center
(Chapter~\ref{ch:drinfeld-center}).

\section{The $E_1$-Koszul duality}
\label{sec:e1-koszul}

\begin{theorem}[$E_1$-chiral Koszul duality, Volume~II]
\label{thm:e1-koszul}
\ClaimStatusProvedElsewhere
For an $E_1$-chiral algebra $\cA$, the bar-cobar adjunction of Volume~I refines to an $E_1$-equivariant adjunction.  On the representation-category level, this gives a monoidal equivalence
\[
  \Rep^{E_1}(A) \simeq \Rep^{E_1}(A^!)^{\mathrm{rev}}
\]
between the monoidal representation category of $A$ and the reverse monoidal category of $A^!$.
\end{theorem}
% AP25: A^! here means the Koszul dual algebra, i.e., the linear dual of
% bar cohomology H*(B(A)). This is NOT the cobar Omega(B(A)) (which
% recovers A itself) nor the Verdier dual D_Ran(B(A)) (which is B(A!)).
% See Vol I, Convention conv:bar-coalgebra-identity.
